\documentclass[a4paper, french, 11pt]{article}

\usepackage[utf8]{inputenc}
\usepackage[T1]{fontenc}
\usepackage[left=3cm, right=3cm, top=1.5cm, bottom=3cm]{geometry}
\usepackage{lmodern}
\usepackage{xspace}
\usepackage{varioref}
\usepackage{float}
\usepackage{xcolor}
\usepackage{siunitx}
\usepackage{amsmath}
\usepackage{amsthm}
\usepackage{tikz}
\usepackage{amssymb}
\usepackage{tabularx}
\usepackage{enumitem}
\usepackage{euler}
\usepackage[many]{tcolorbox}
\usepackage{babel}
\usepackage{listings}
\usepackage{biblatex}
\usepackage{csquotes}% Recommended

% LISTING CONFIG
\newtoggle{InString}{}% Keep track of if we are within a string
\togglefalse{InString}% Assume not initally in string
\newcommand*{\ColorIfNotInString}[1]{\iftoggle{InString}{#1}{\color{green!60!black}#1}}%
\newcommand*{\ProcessQuote}[1]{#1\iftoggle{InString}{\global\togglefalse{InString}}{\global\toggletrue{InString}}}%
\lstset{
  literate=
  {0}{{{\ColorIfNotInString{0}}}}1
  {1}{{{\ColorIfNotInString{1}}}}1
  {2}{{{\ColorIfNotInString{2}}}}1
  {3}{{{\ColorIfNotInString{3}}}}1
  {4}{{{\ColorIfNotInString{4}}}}1
  {5}{{{\ColorIfNotInString{5}}}}1
  {6}{{{\ColorIfNotInString{6}}}}1
  {7}{{{\ColorIfNotInString{7}}}}1
  {8}{{{\ColorIfNotInString{8}}}}1
  {9}{{{\ColorIfNotInString{9}}}}1
}
\lstset{
  literate=
  {á}{{\'a}}1 {é}{{\'e}}1 {í}{{\'i}}1 {ó}{{\'o}}1 {ú}{{\'u}}1
  {Á}{{\'A}}1 {É}{{\'E}}1 {Í}{{\'I}}1 {Ó}{{\'O}}1 {Ú}{{\'U}}1
  {à}{{\`a}}1 {è}{{\`e}}1 {ì}{{\`i}}1 {ò}{{\`o}}1 {ù}{{\`u}}1
  {À}{{\`A}}1 {È}{{\'E}}1 {Ì}{{\`I}}1 {Ò}{{\`O}}1 {Ù}{{\`U}}1
  {ä}{{\"a}}1 {ë}{{\"e}}1 {ï}{{\"i}}1 {ö}{{\"o}}1 {ü}{{\"u}}1
  {Ä}{{\"A}}1 {Ë}{{\"E}}1 {Ï}{{\"I}}1 {Ö}{{\"O}}1 {Ü}{{\"U}}1
  {â}{{\^a}}1 {ê}{{\^e}}1 {î}{{\^i}}1 {ô}{{\^o}}1 {û}{{\^u}}1
  {Â}{{\^A}}1 {Ê}{{\^E}}1 {Î}{{\^I}}1 {Ô}{{\^O}}1 {Û}{{\^U}}1
  {œ}{{\oe}}1 {Œ}{{\OE}}1 {æ}{{\ae}}1 {Æ}{{\AE}}1 {ß}{{\ss}}1
  {ű}{{\H{u}}}1 {Ű}{{\H{U}}}1 {ő}{{\H{o}}}1 {Ő}{{\H{O}}}1
  {ç}{{\c c}}1 {Ç}{{\c C}}1 {ø}{{\o}}1 {å}{{\r a}}1 {Å}{{\r A}}1
  {€}{{\EUR}}1 {£}{{\pounds}}1
}
\definecolor{codegray}{rgb}{0.5,0.5,0.5}
\definecolor{backcolour}{rgb}{0.95,0.95,0.92}
\lstdefinestyle{myStyle}{
  backgroundcolor=\color{backcolour},
  commentstyle=\color{blue!50!white},
  keywordstyle=\color{orange},
  numberstyle=\small\selectfont\color{codegray},
  stringstyle=\color{red!80!orange},
  basicstyle=\ttfamily\small\selectfont,
  breakatwhitespace=false,
  breaklines=true,
  captionpos=b,
  keepspaces=true,
  numbers=left,
  numbersep=5pt,
  showspaces=false,
  showstringspaces=false,
  showtabs=false,
  tabsize=4,
  xleftmargin=.025\textwidth,
  xrightmargin=.025\textwidth
}

\addbibresource{refs.bib}% Syntax for version >= 1.2

\newtheoremstyle{break}
{}{}
{}{}
{\bfseries}{}
{\newline}{}
\theoremstyle{break}
\newtheorem{Question}{Question}


\newcommand{\po}{\left(}
\newcommand{\pf}{\right)}

\date{}
\title{Gestion des scopes pour Python}
\author{Martin \textsc{Andrieux}}

\begin{document}
\maketitle

\section{Gestion des scopes dans la thèse de Raphaël Monat}
Dans la thèse, il y a un unique environnement \(\mathcal E \overset{def}{=} \mathsf{Id} \rightharpoonup \mathsf{Addr} \cup \mathtt{LocalUndef}\).
Les mots clés \texttt{global} et \texttt{nonlocal} ne sont pas supportés.
Une analyse est faite pendant le parsing pour détecter les variables non déclarées.

Nous souhaitons gérer différemment les portées des variables (scopes), afin de pouvoir utiliser \texttt{global} et \texttt{nonlocal}.
Nous expliquons notre approche dans la section \vref{pyskel}, après une brève explication des scopes en Python.

\section{Scopes en Python}
Nous détaillons dans cette section les principaux aspects de la gestion des scopes en python à travers divers exemples.

Les variables et fonctions définies en dehors de toute fonction/classe sont dans le \emph{scope global}.

Dans le corps d'une fonction, le scope de contient que les variables locales, essayer de modifier une variable globale mènera à la création d'une nouvelle variable locale à la fonction.
Par exemple, dans le programme de la figure \vref{fresh_scope}, la variable globale \texttt{x} est toujours égale à \texttt{42} après l'appel de \texttt{f}.

\begin{figure}[h]
  \centering
  \begin{minipage}{0.5\linewidth}
\begin{lstlisting}[style=myStyle, language=python]
  x = 42 # Global x

  def f():
    x = 0 # x is local to f

  f()
  print(x) # => 42\end{lstlisting}
\end{minipage}
\caption{\label{fresh_scope}La valeur de \texttt{x} (global) reste inchangée après l'appel de \texttt{f}}
\end{figure}

Il est néanmoins possible d'accéder à des variables globales depuis le corps des fonctions grâce à l'instruction \texttt{global}, comme démontré figure \vref{global_keyword}.

\begin{figure}[h]
  \centering
  \begin{minipage}{0.7\linewidth}
\begin{lstlisting}[style=myStyle, language=python]
  x = 42 # Global x

  def f():
    # Says "The x i'm using is the global x"
    global x
    x = 0

  f()
  print(x) => 0
\end{lstlisting}
  \end{minipage}
  \caption{\label{global_keyword}Utilisation du mot clé \texttt{global}}
\end{figure}

Finalement, le mot clé \texttt{nonlocal} permet d'accéder au scope \og au dessus\fg{} en cas de fonctions imbriquées, voir figure \vref{nonlocal_keyword}.

\begin{figure}[h]
  \centering
  \begin{minipage}{0.7\linewidth}
\begin{lstlisting}[style=myStyle, language=python, morekeywords={nonlocal}]
def g():
  x = 42 # x is not global (it's local to g)

  def f():
    # global x will fail
    # We use the x of g
    nonlocal x
    x = 0

  f()
  print(x)

g() => 0
\end{lstlisting}
  \end{minipage}
  \caption{\label{nonlocal_keyword}Utilisation du mot clé \texttt{nonlocal}}
\end{figure}

Notez que \texttt{nonlocal} peut remonter plusieurs fonction pour trouver la variable déclarée.

Sachant tout cela, il est possible de partager de la mémoire entre fonctions, comme illustré figure \vref{shared_memory}.
Les deux fonctions \texttt{init} et \texttt{incr} modifient la même zone mémoire, à savoir la variable \texttt{x} locale à \texttt{g}.
\begin{figure}[h]
  \centering
  \begin{minipage}{0.7\linewidth}
\begin{lstlisting}[style=myStyle, language=python, morekeywords={nonlocal}]
def g():
  x = 42

  def init():
    nonlocal x
    x = 0

  def incr():
    nonlocal x
    x += 1

  return (init, get)
\end{lstlisting}
  \end{minipage}
  \caption{\label{shared_memory}Partage de mémoire entre sous-fonctions}
\end{figure}

\section{\label{pyskel}Gestion des scopes dans pyskel}
\subsection{Types}

Nous découpons l'environnement initial (celui de la thèse de Raphaël Monat), en un environnement et un scope.

\begin{itemize}
  \item L'environnement associe les noms de variables à une portée (\texttt{Local}, \texttt{Nonlocal} ou \texttt{Global})
  \item Le scope associe les identifiant à une adresse dans le tas. Il contient également le scope supérieur s'il existe.
\end{itemize}

Ainsi, pour retrouver l'adresse associée à un identifiant, nous commençons par regarder le scope de ce dernier (contenu dans l'environnement).
\begin{itemize}
  \item S'il est local, nous regardons dans le scope local
  \item s'il est global, nous regardons dans le scope global
  \item s'il est nonlocal, nous regardons récursivement dans les scopes au dessus.
\end{itemize}

Nous définissons donc les types suivants:
\begin{lstlisting}[style=myStyle, language=caml]
type var_scope =
  | Local
  | NonLocal
  | Global

(* String map *)
type map<_>
(* Heaps map addr to values *)
type heap<_>
type heap_addr<_>

type scope_id = heap_addr<scope>

(* Read-only *)
type env =
  | Globl
  | InFun (map<var_scope>, scope_id)

type scope_heap = heap<scope>

type scope =
  ( scope_up: maybe<scope_id>
  , mapping: map<maybe<addr>>
  )
\end{lstlisting}

En plus de celà, nous définissions un scope global. Nous utilisons directement des \texttt{addr} car il n'y a jamais d'initialisation à \texttt{LocalUndef} (représenté par \texttt{Nothing}) dans le scope global.
\begin{lstlisting}[style=myStyle, language=caml, firstnumber=17]
type global_scope = map<addr>
\end{lstlisting}

L'état interne de notre interpréteur est alors défini comme suit:
\begin{lstlisting}[style=myStyle, language=caml]
(* Read only state *)
type r = (local_env: env, ...)

(* Read Write state *)
type s =
  ( global_scope: global_scope
  , scope_heap: scope_heap
  , ...
  )
\end{lstlisting}

Les mappings vers les adresses ne peuvent pas être en lecture seule car ils sont modifiés lors des assignations.

L'utilisation d'un tas de scope est nécessaire au partage de mémoire.
En effet, dans l'exemple de la figure \vref{shared_memory}, les scopes de \texttt{init} et \texttt{incr} possèdent le même \texttt{scope\_up}, et font donc effectivement référence à la même zone mémoire.
Nous allouons donc les scopes, comme nous allouons des valeurs python (strings entiers, fonctions, \dots).

\subsection{Recherche d'une valeur}
La fonction \texttt{readEnv : id -> m<addr>} permet d'obtenir l'adresse correspondant à l'identifiant passé en paramètre.
Une version simplifiée est donnée figure \vref{readEnv}.

\begin{figure}[h]
  \centering
  \begin{minipage}{\textwidth}
\begin{lstlisting}[style=myStyle, language=caml]
val readEnv (id:id) : m<addr> =
  let r =@ ask in
  match r.local_env with
  | Globl ->
  (* We are in the global env, we just look for the variable in the global map *)
    find_global id
  | InFun (e, scope_id) ->
  (* We are evaluating a function, we refer to the local environment *)
    let scope_var = map_lookup<var_scope> (id, e) in
    match scope_var with
    | Nothing ->
    (* Not in env, but it's a read, so we lookup upper scopes *)
      let ma =@ find_nonlocal (id, scope_id) in
      match ma with
      | Nothing -> find_global id
      | Just (Just a) -> ret a
      | Just Nothing -> fail "Out of scope"
      end
    (* Look in the right scope *)
    | Just Local ->
        let (Just a) =@ find_local (id, scope_id) in
        match a with
        | Just a -> ret a
        (* LocalUndef *)
        | Nothing -> fail "id is bound to LocalUndef"
        end
    | Just NonLocal ->
        let Just (Just a) =@ find_nonlocal (id, scope_id) in
        ret a
    | Just Global -> find_global id
    end
  end
\end{lstlisting}
  \end{minipage}
  \caption{\label{readEnv}Fonction \texttt{readEnv}}
\end{figure}

La fonction \texttt{find\_local} regarde si l'identifiant est dans le scope local, et renvoie l'adresse associée. \texttt{find\_nonlocal} fait de même de manière récursive dans les sous environnements.

La fonction \texttt{writeEnv} est très similaire à \texttt{readEnv}, la tâche difficile étant de trouver la bonne occurrence de l'identifiant dans le scope.

\subsection{Définition de fonction}
Lors de la définition d'une fonction, nous devons capturer l'identifiant du scope courant, contenu dans l'environnement en lecture seule.
\begin{lstlisting}[style=myStyle, language=caml]
  | SFunctionDef (name, args, body) ->
    let bodyLocals = locals body in
    let e =@ ask in
    let scope_up =
      match e.local_env with
      | Globl -> Nothing<scope_id>
      | InFun(_, sid) -> Just<scope_id> sid
      end in
    let f = Fun(name, args, bodyLocals, body, scope_up) in
    let addr =@ alloc (f, Locked) in
    writeEnv (name, addr)
\end{lstlisting}

\subsection{Appel de fonction}
L'appel de fonction nécessite d'allouer un nouveau scope (le scope courant dans le corps de la fonction).
Ce scope aura comme \og supérieur\fg{} le scope contenu dans l'objet \texttt{Fun}.

Il suffit ensuite de lancer l'évaluation de la fonction en modifiant localement l'état en lecture seule. Les annotations de types ont été retirées pour plus de lisibilité.

\begin{lstlisting}[style=myStyle, language=caml]
  | Call(e_c, args) ->
    let r =@ ask in
    (* Eval the function *)
    let at_c =@ eval_expr e_c in
    let (Fun(_, a, loc, body, scope_up), _) =@ read at_c in

    (* Eval and set the arguments *)
    let at_args =@ mapM (eval_expr, args) in
    let args_with_values = zip (a, at_args) in

    (* Set the local variables to LocalUndef *)
    let locals_undef =
      let f ((m, s):(map, id)) = map_insert (s, Nothing, m) in
      foldl (f , map_empty , loc) in
    let locals_with_args =
      let f ((m, (s, a)):(map, (id, addr))) = map_insert (s, Just a, m) in
      foldl (f , locals_undef , args_with_values) in

    (* Scope Allocation *)
    let inner_scope = (scope_up = scope_up, mapping = locals_with_args) in
    let inner_scope_id =@ allocScope inner_scope in

    let non_locals = nonlocals body in
    let global_vars = globals body in
    let all_locals = append (loc, a) in
    let fun_env =
      let is_local ((m, s):(map, id)) = map_insert (s, Local, m) in
      let is_nonlocal ((m, s):(map, id)) = map_insert (s, NonLocal, m) in
      let is_global ((m, s):(map, id)) = map_insert (s, Global, m) in
      let m = foldl (is_local , map_empty , all_locals) in
      let m' = foldl (is_nonlocal, m, non_locals) in
      foldl (is_global, m', global_vars) in

    let set_locals = \_:r ->
      let local_env' = InFun (fun_env, inner_scope_id) in
      r <- (local_env = local_env')
    in

    let bodyEvaluation = (
      eval_stmts body ;@
      return r.builtins.none
    ) in

    let localBodyEval = local (set_locals, bodyEvaluation) in

    let catchReturn = \f:flag ->
      match f with
      | Ret a -> return a
      | f -> setFlag f
      end in

    (* Eval the body of the function *)
    catchFlag (localBodyEval, catchReturn)
  end
\end{lstlisting}

\end{document}
